\documentclass[11pt]{article}
\usepackage{amsmath,amssymb}
\title{A logarithmic signed magnetic bound with the original mixed boundaries}\author{}\date{}
\begin{document}\maketitle
\paragraph{Objective and scope.}
The constant flux spin and magnetic reduction calculation reduces the ORIGINAL complete constant-flux quadratic response to
four scalar magnetic edge determinants, twice the scalar ghost determinant,
and already bounded spin/potential/neutral terms. Its spin trace estimate
and boundary-normal diagonal check are used. Here an elementary killed-walk
argument controls the remaining scalar determinants with their actual
mixed boundary conditions. This closes the leading quadratic estimate for
ONE physical direction. The nonlinear original-law remainder (5) of the signed gauge quotient and unsummed current calculation
is not estimated by this argument.

\paragraph{1. The literal scalar connections.}
Put $f=R^{-2}$. In the conventions of the constant flux spin and magnetic reduction calculation, the forward scalar hopping
phase on the component-$\mu$ edge grid in direction $\nu$ is
\[
 a^{(\mu)}_\nu(x)=f\left[
  \mathbf1_{\nu=1}(2x_0+\mathbf1_{\mu=0})
       +\mathbf1_{\nu=0}\mathbf1_{\mu=1}\right].                    \tag{1}
\]
This follows directly from the same-component ordered Wilson insertions
and the endpoint gauge term. In the direction of the link component the
hopping comes from the gauge term; in a transverse direction it comes
from its single incident plaquette. The ghost phase is
$a^{(g)}_\nu(x)=2f x_0\mathbf1_{\nu=1}$.
The extra terms in (1) are exact gradients of
$f(\mathbf1_{\mu=0}x_1+\mathbf1_{\mu=1}x_0)$.
Hence every scalar operator is unitarily gauge equivalent to the SAME
constant-flux connection $a_1=2fx_0$, on its OWN finite grid.
This does not equate their boundary conditions. Each edge grid has $R$
Neumann sites in direction $\mu$, and $R-1$ Dirichlet interior sites in
each other direction; the ghost has four Dirichlet directions.
The original $R=3$ full-Hessian derivative check verifies every elementary
square curvature $2f$ in plane01 and zero otherwise, error below $6\,10^{-17}$.

\paragraph{2. A scalar mixed-boundary lemma.}
Let $L_\epsilon$ be the nearest-neighbor magnetic difference matrix on any
of these grids, with hopping phases $\epsilon a$, flat original diagonal,
and connection gauge equivalent to $a_1=b x_0$, $b=2f$.
Then with an absolute constant independent of $R$,
\[
 \boxed{0\le (\log\det L_\epsilon)''\big|_0
            \le C |\Omega|b^2\log(2+R)\le C\log(2+R).}             \tag{2}
\]
The first bound's factor $|\Omega|b^2$ is retained: $|\Omega|\le R^4$
and $b^2=4R^{-4}$ only in this specific physical direction.

Here is a direct proof, including the reflection issue. Realize $L_0$
as the negative generator of independent continuous-time coordinate walks,
with rate one to each nearest neighbor. A Neumann endpoint suppresses the
outward jump; a Dirichlet endpoint kills the walk on the outward jump.
Write $\tau$ for the killing time. At least three coordinates are
Dirichlet, so the original operator is strictly positive.

Fix a starting point $x$. Among coordinates zero and one at least ONE is
Dirichlet, since there is at most one Neumann coordinate in the entire
grid. Choose it as $y$, and call the other $z$. A gauge change puts the
connection in the form
\[
 a_y=\pm b(z-z(x)),\qquad a_\nu=0\quad(\nu\ne y).                  \tag{3}
\]
The subtraction of $z(x)$ is a constant-gradient gauge change. It depends
on the chosen start, which is harmless for a closed-loop trace.
Extend the coordinate processes beyond the killing time, and stop their
current at $\tau$. Because $y$ is Dirichlet, its two jumps have equal
rates until killing, INCLUDING the possible killing jump. Thus
\[
 M_t=\sum_{\substack{y\text{-jumps at }s\le t\wedge\tau}}
             \pm b(Z_{s-}-z(x))\,\Delta Y_s
\]
is a stopped martingale. Its predictable quadratic variation is
$2b^2\int_0^{t\wedge\tau}(Z_s-z(x))^2ds$.
If $z$ is Dirichlet, use the unrestricted integer walk up to killing.
If $z$ is Neumann, its reflected finite-interval walk is a 1-Lipschitz
fold of the unrestricted integer walk (the suppressed endpoint jump
becomes a self jump under the fold). Consequently, for either case,
$E(Z_s-z(x))^2\le2s$ for the extended process. Dropping the stopping
indicator in the nonnegative integral proves
\[
 E_x M_t^2\le 2b^2t^2,
 \qquad E_x[M_t^2;\tau>t]\le2b^2t^2.                              \tag{4}
\]
No reflected-coordinate drift estimate is assumed: the current is
integrated only against a DIRICHLET coordinate. This choice is what uses
the actual mixed boundary structure.

\paragraph{3. Close the loop before bounding its current.}
Let $p_t$ be the original zero-field killed transition kernel. The
one-dimensional sine/cosine product expansions give
\[
 \sup_{u,v}p_t(u,v)\le C(1+t)^{-2}e^{-ct/R^2}.                     \tag{5}
\]
For completeness each Dirichlet factor is bounded by
$C(1+t)^{-1/2}e^{-ct/R^2}$, and the one possible Neumann factor by
$C[(1+t)^{-1/2}+R^{-1}]$. For $t\le R^2$ the first term controls both;
for $t>R^2$ the extra factor $\sqrt{t}/R$ is absorbed by a smaller
constant in the exponential supplied by the three Dirichlet factors.
The same estimate holds with four Dirichlet factors.

The magnetic exponential expansion is the jump-path expansion of the
same walk, each surviving path weighted by $e^{i\epsilon J}$, where $J$
is its oriented current. In particular
\[
 -\partial_\epsilon^2\operatorname{Tr}e^{-tL_\epsilon}\big|_0
 =\sum_x E_x[J_t^2;X_t=x,\tau>t]=:\mathcal A(t).                  \tag{6}
\]
Split a closed path into two halves of length $s=t/2$ and reverse the
second. Reversibility of the killed real kernel makes both halves paths
from $x$ to the same midpoint; their weights multiply. The total current
is $J_1-J_2$. With the SAME start-centered gauge (3),
$(J_1-J_2)^2\le2J_1^2+2J_2^2$. For each half, sum its squared current
first and bound the other half by $\sup p_s$. Equations(4)--(5) give
\[
 \mathcal A(t)\le4|\Omega|(\sup p_s)\sup_xE_x[M_s^2;\tau>s]
 \le C|\Omega|b^2\frac{t^2}{(1+t)^2}e^{-ct/R^2}.                  \tag{7}
\]
This keeps the four-dimensional return probability. Using an unconditioned
area bound followed by Cauchy--Schwarz would lose that power.

The finite positive-matrix logarithm identity, differentiated twice, gives
\[
 (\log\det L_\epsilon)''\big|_0
 =\int_0^\infty\mathcal A(t)\,\frac{dt}{t}.
\]
All terms are nonnegative. Formula(7) proves convergence at zero and
infinity and gives
$\int_0^\infty t(1+t)^{-2}e^{-ct/R^2}dt\le C\log(2+R)$,
which proves (2). The finite-dimensional differentiation can equivalently
be done first with finite time cutoffs and then passed to the limit using
(7); no infinite-volume limit or unproved bridge estimate is used.

\paragraph{4. Consequence for the original signed quadratic response.}
Apply (2) to each of the four ORIGINAL scalar edge operators and to the
original ghost operator. The signed combination of the constant flux spin and magnetic reduction calculation therefore obeys
$|\mathcal M_R|\le C\log(2+R)$. Its exact spin remainder costs at most the
same logarithm, while potential and neutral terms cost $C$. Consequently
for the actual constant-flux minimum of the complete cube abelian signed incidence calculation,
\[
 \boxed{|F_{{\rm quad},\zeta\zeta}|
       \le {C\log(2+R)\over B}
       \le {C\log(2+R)\over B}\,E_R,\qquad
       \epsilon=\zeta/\sqrt B,\quad E_R={(R+1)^2\over R^2}.}         \tag{8}
\]
This is a signed leading quadratic physical-energy bound with the original
relative boundary conditions and full gauge normalization. No replacement
of a global holonomy law by free rotors, or deletion of a live determinant,
occurs. Its constant does not determine the renormalized coupling coefficient.
The original compact nonlinear integral still differs by the exact
second/fourth-moment and current terms in (5) of the signed gauge quotient and unsummed current calculation. Controlling those by
this same energy, and extending from one flux direction to general admitted
nonflat records under composition, are indispensable open steps. Thus (8)
is not a full nonlinear block estimate, nor S94 or a physical-scale induction.
\end{document}
